\chapter{Elicitation and Expert Data}
\label{chapter:elicitation}

\section{The Expert Panel}
\label{sec:expert_panel}

The seven-member panel was deliberately composed to combine three forms of substantive domain knowledge: institutional research, market-facing commercial operations, and lived consumer experience of the Dutch housing crisis. This stratification is important because the target variables span construction, transactions, and rents. No single professional viewpoint covers that full system. The panel therefore combines formal knowledge of housing policy and management, daily exposure to property and rental markets, and direct experience of localized scarcity.

Dr.~Sake Zijlstra is an Assistant Professor and lecturer in Housing Management at the Faculty of Architecture and the Built Environment, Delft University of Technology. His contribution is grounded in institutional research on housing systems, management, and policy. Esther den Dikken works in real-estate operations at Key Vastgoed Beheer, a property-management and rental agency in the Netherlands. Her daily market exposure provides frontline insight into free-sector rental transactions, market pressure, and commercial vacancy dynamics. Marijke de Vries is an MSc Architecture and Urbanism student at TU Delft whose academic work focuses on Dutch housing supply, spatial planning, and affordability policy.

Mario Gonzalez, Gabriel Rodriguez, Carlos Hita, and Paul Casamian are university students who have lived in the Netherlands for several years. Each has navigated the acute student-housing shortage directly and therefore contributes lived knowledge of local search costs, scarcity, and supply-demand friction. Their practical experience is supplemented by independent academic study of macroeconomic housing indicators. In the Classical Model, substantive knowledge need not be identical across participants; the calibration variables test whether these different forms of knowledge translate into statistically accurate and informative assessments.

\begin{table}[H]
    \centering
    \small
    \begin{tabularx}{\textwidth}{p{3.5cm}X}
        \toprule
        \textbf{Panel member} & \textbf{Qualification and domain contribution} \\
        \midrule
        Dr.~Sake Zijlstra & Assistant Professor/lecturer in Housing Management, TU Delft; housing systems, management, and policy \\
        Esther den Dikken & Real-estate professional at Key Vastgoed Beheer; property management, free-sector rentals, market pressure, and vacancy \\
        Marijke de Vries & MSc Architecture and Urbanism student, TU Delft; housing supply, spatial planning, and affordability policy \\
        Mario Gonzalez & University student and long-term Netherlands resident; lived experience of student-housing scarcity and independent study of housing indicators \\
        Gabriel Rodriguez & University student and long-term Netherlands resident; lived experience of student-housing scarcity and independent study of housing indicators \\
        Carlos Hita & University student and long-term Netherlands resident; lived experience of student-housing scarcity and independent study of housing indicators \\
        Paul Casamian & University student and long-term Netherlands resident; lived experience of student-housing scarcity and independent study of housing indicators \\
        \bottomrule
    \end{tabularx}
    \caption{Panel qualifications and complementary forms of substantive domain knowledge.}
    \label{tab:expert_profiles}
\end{table}

Accountability and anonymity are managed at different levels. Names and qualifications are reported here so that readers can evaluate whether the panel possesses relevant expertise and so that the researcher remains accountable for recruitment. Individual quantitative assessments are nevertheless scrubbed of names in Chapters~4 and~5, where participants are referred to only as Expert 1 through Expert 7. The mapping between names and analytical identifiers is not published. This prevents professional reputation, hierarchy, or peer pressure from shaping the interpretation of individual scores while preserving an auditable internal record. The raw export contains personal data and is not included in the public repository; only anonymized derived results and code are retained.

\section{Elicitation Format and Dry Run}
\label{sec:elicitation_format}

For each calibration and target variable, experts provided the 5th, 50th, and 95th percentiles of their subjective distribution. The 5th and 95th percentiles form a central 90\% uncertainty interval; the 50th percentile is the median. Participants were instructed that 5\% of their probability should lie below the lower value and 5\% above the upper value.

The dry run identified two problems. First, participants tended to treat the median as a deterministic ``best guess'' and then place narrow bounds around it. Second, some mixed percentages, proportions, billions of euros, and absolute amounts. The final instrument therefore displayed units beside every input, visually separated percentage questions from count and currency questions, enforced $q_{0.05}<q_{0.50}<q_{0.95}$, and replaced a central point-estimate example that encouraged anchoring with a neutral interval illustration. Participants were asked to consider plausible low and high outcomes before entering the median.

These changes reduced avoidable ambiguity but did not eliminate it completely. Two final response patterns still suggest scale misinterpretation. They were retained rather than corrected after elicitation, because post-hoc adjustment would compromise empirical control; their effect is instead reflected in the calibration scores.

\section{Calibration Questions}
\label{sec:calibration_questions}

Ten seed variables were chosen from the same broad domain as the target questions. They cover transaction prices, residential mortgage debt, population growth, housing stock, rental tenure, existing-home transactions, one-person households, international student enrolment, rent-to-income ratios, and homelessness. Their realizations were fixed before scoring and concealed from participants during elicitation.

The realizations were obtained from traceable Dutch sources, principally CBS, the BAG register, De Nederlandsche Bank, the NVM quarterly market report, and the Dutch Housing Survey. Appendix~\ref{appendix:elicitation_protocol} gives the complete wording, units, realizations, and source for every seed and target variable. Using multiple source types broadens domain coverage, although correlation between housing indicators remains a limitation discussed in Chapter~\ref{chapter:conclusion_discussion}.
